\documentclass[11pt]{article}
\usepackage[a4paper,margin=1in]{geometry}
\usepackage{amsmath,amssymb,amsthm,booktabs,microtype}
\usepackage{cleveref}
\newtheorem{theorem}{Theorem}[section]
\newtheorem{corollary}[theorem]{Corollary}
\newtheorem{proposition}[theorem]{Proposition}
\theoremstyle{remark}
\newtheorem{remark}[theorem]{Remark}
\newcommand{\Tr}{\operatorname{Tr}}
\newcommand{\1}{\mathbf 1}
\newcommand{\cA}{\mathcal A}
\newcommand{\cB}{\mathcal B}
\newcommand{\cP}{\mathcal P}
\newcommand{\cR}{\mathcal R}
\newcommand{\cS}{\mathcal S}
\newcommand{\cT}{\mathcal T}
\title{Parity-Renormalized First-Alias Packets:\\Exact Signs, Clock Phase, and the Neighboring-Shell Interface}
\author{Bin Wang}
\date{July 2026}
\begin{document}
\maketitle

\begin{abstract}
The first counterloop alias and the square-root parity leakage occur at the
same weighted order, so their signs and normalization cannot be suppressed.
We restore the Hardy scaling hidden in the earlier first-alias ledger and
split the coefficient defect exactly into a raw-trace packet, a signed parity
packet, a finite-radius radial correction, and the roots-of-unity alias
impulse.  At the even first alias the parity correction is positive, whereas
the counterloop defect is subtracted from the total residual.  The parity
packet has a uniform moving-order expansion with an explicit quadratic
remainder.  On a fixed clock phase $\eta$, its leading ratio to the alias
defect is $C_*C_M\lambda^\eta$, so the unique symbolic balance phase is
$-\log(C_*C_M)/\log\lambda$.  If $C_*C_M\lambda<1$, this gives a strict
scalar-only obstruction throughout the usual floor/ceil/nearest phase window
$|\eta|\le1$.  The archived decimal evaluation supports that hypothesis but
is not an interval certificate.  This conditional negative result does not
exclude cancellation by the raw physical boundary packet and neighboring
shell.  We finish with a phase- and coordinate-preserving interface for that
calculation.  No joint full-trace replacement is proved.
\end{abstract}

\section{Normalization and four packet types}

Put
\begin{equation}
 \lambda=1.678573510428322\ldots,
 \qquad r_H=0.85,
 \qquad R=1.4,
 \qquad \beta=(r_H\sqrt\lambda)^{-1}.
 \label{eq:constants}
\end{equation}
The number $r_H$ is the chosen Hardy scaling radius, not a dynamical
invariant.  Let $P_n$ be the deterministic physical flat trace and let
$K_\sigma$ be the row-normalized noisy Markov operator.  The peripheral
parity eigenvalue is written
\begin{equation}
 \lambda_-(\sigma)=-(1-\delta_\sigma),
 \qquad
 \delta_\sigma=C_*\sqrt\sigma+o(\sqrt\sigma),
 \qquad
 C_*=0.105258535936908\ldots .
 \label{eq:parity-gap}
\end{equation}
The decimal evaluates an archived positive integral; positivity and the
square-root law, rather than the rounded decimal, are the analytic inputs.

We use explicit Hardy superscripts to avoid mixing two archived conventions:
\begin{align}
 c^H_{\sigma,n}
 &=r_H^{-n}\{\Tr K_\sigma^n-1-\lambda_-(\sigma)^n\},
 \label{eq:noisy-hardy}\\
 c^H_n
 &=r_H^{-n}\{P_n-1-(-1)^n\}.
 \label{eq:det-hardy}
\end{align}
Define the raw-trace and parity packets
\begin{align}
 \cT_{\sigma,n}
 &=r_H^{-n}\{\Tr K_\sigma^n-P_n\},
 \label{eq:raw-packet}\\
 \cP_{\sigma,n}
 &=r_H^{-n}\{(-1)^n-\lambda_-(\sigma)^n\}
 =(-1)^nr_H^{-n}\{1-(1-\delta_\sigma)^n\}.
 \label{eq:parity-packet}
\end{align}
Then, identically,
\begin{equation}
 c^H_{\sigma,n}-c^H_n=\cT_{\sigma,n}+\cP_{\sigma,n}.
 \label{eq:bulk-split}
\end{equation}

For $k\ge2$, let $M_k$ be the archived boundary-cycle multiplier and put
\begin{equation}
 \rho_k=|M_k|^{-1/k},
 \qquad
 \beta_k=\frac{\sqrt{\rho_k}}{r_H}
 =\frac{|M_k|^{-1/(2k)}}{r_H}.
 \label{eq:beta-type}
\end{equation}
The Hardy-scaled finite counterloop is
\begin{equation}
 s_{k,n}=\beta_k^n
 \bigl(2k\1_{2k\mid n}-1-(-1)^n\bigr),
 \qquad \beta_k\longrightarrow\beta.
 \label{eq:counterloop}
\end{equation}
Rename the limiting pole moment and numerator anchor to avoid collisions with
older uses of $p_n$ and $a_n$:
\begin{equation}
 p_n^{\rm pole}=-2\1_{2\mid n}\beta^n,
 \qquad
 a_n^{\rm num}=c^H_n-p_n^{\rm pole}.
 \label{eq:pole-anchor}
\end{equation}

\section{Exact alias and parity algebra}

\begin{theorem}[Radial-plus-alias decomposition]
\label{thm:counterloop-split}
For every $k\ge2$ and $n\ge1$, the counterloop defect is exactly
\begin{equation}
 \boxed{
 \cA_{k,n}:=s_{k,n}-p_n^{\rm pole}
 =2\1_{2\mid n}(\beta^n-\beta_k^n)
 +2k\beta_k^n\1_{2k\mid n}.}
 \label{eq:counterloop-split}
\end{equation}
The first term is finite-radius radial drift and the second is the exact
roots-of-unity alias impulse.  In particular, $\cA_{k,n}=0$ for every odd
$n$, while
\begin{equation}
 \cA_{k,2k}=(2k-2)\beta_k^{2k}+2\beta^{2k}>0.
 \label{eq:first-alias-defect}
\end{equation}
\end{theorem}

\begin{proof}
Substitute \eqref{eq:counterloop} and \eqref{eq:pole-anchor}.  If $n$ is
odd, both $1+(-1)^n$ and the two divisibility indicators vanish.  If $n$ is
even, the non-alias part is $-2\beta_k^n+2\beta^n$, and an order divisible by
$2k$ receives the additional impulse $2k\beta_k^n$.  Setting $n=2k$ gives
\eqref{eq:first-alias-defect}.
\end{proof}

The sign ledger is now unambiguous:
\begin{center}\small
\begin{tabular}{lccc}
\toprule
order & parity packet $\cP_{\sigma,n}$
& counterloop defect $\cA_{k,n}$ & signs in the residual\\
\midrule
odd & negative & $0$ & $+\cP$\\
even, $2k\nmid n$ & positive
& $2(\beta^n-\beta_k^n)$ & $+\cP-\cA$\\
first alias $n=2k$ & positive
& $(2k-2)\beta_k^{2k}+2\beta^{2k}$ & $+\cP-\cA$\\
\bottomrule
\end{tabular}
\end{center}
Only the parity sign is asserted in the middle row; the radial difference
need not be assigned a sign without a finite-$k$ radius inequality.

\begin{theorem}[Parity-renormalized first-alias packet identity]
\label{thm:main-identity}
Let
\begin{equation}
 e_{\sigma,k,n}=c^H_{\sigma,n}-s_{k,n}-a_n^{\rm num}.
 \label{eq:error}
\end{equation}
Then at the first alias,
\begin{align}
 e_{\sigma,k,2k}
 &=\cT_{\sigma,2k}+\cP_{\sigma,2k}-\cA_{k,2k}
 \label{eq:packet-form}\\
 &=r_H^{-2k}\{\Tr K_\sigma^{2k}-P_{2k}
     +1-\lambda_-(\sigma)^{2k}\}
 \notag\\[-0.2em]
 &\hspace{5em}
 -\{(2k-2)\beta_k^{2k}+2\beta^{2k}\}.
 \label{eq:expanded-form}
\end{align}
In particular, the scalar parity contribution in \eqref{eq:expanded-form}
is
\begin{equation}
 r_H^{-2k}\{1-\lambda_-(\sigma)^{2k}\}
 =r_H^{-2k}\{1-(1-\delta_\sigma)^{2k}\}>0.
 \label{eq:positive-even-parity}
\end{equation}
\end{theorem}

\begin{proof}
Using $a_n^{\rm num}=c^H_n-p_n^{\rm pole}$ in \eqref{eq:error} gives
\[
 e_{\sigma,k,n}=(c^H_{\sigma,n}-c^H_n)
 -(s_{k,n}-p_n^{\rm pole}).
\]
Apply \eqref{eq:bulk-split}, \cref{thm:counterloop-split}, and set $n=2k$.
Equation \eqref{eq:positive-even-parity} follows because
$0<1-\delta_\sigma<1$ for sufficiently small noise.
\end{proof}

This theorem is an exact coefficient identity.  It does not approximate
$\cT_{\sigma,2k}$ by a local Gaussian law and therefore does not assume the
missing trace-observation theorem.

\section{Uniform parity packet on the first-alias clock}

\begin{theorem}[Uniform scalar parity expansion]
\label{thm:parity-expansion}
For $0\le\delta<1$ and every $n\ge1$,
\begin{equation}
 \left|
 \cP_{\sigma,n}-(-1)^n n\delta r_H^{-n}
 \right|
 \le \frac{n(n-1)}2\delta^2r_H^{-n},
 \label{eq:uniform-remainder}
\end{equation}
where the left side uses \eqref{eq:parity-packet} with
$\delta_\sigma=\delta$.  Consequently, if
\begin{equation}
 k_\sigma=\frac{\log(1/\sigma)}{2\log\lambda}+O(1),
 \label{eq:natural-clock}
\end{equation}
then
\begin{equation}
 \cP_{\sigma,2k_\sigma}
 =2k_\sigma C_*\sqrt\sigma\,r_H^{-2k_\sigma}
 \{1+o(1)\}.
 \label{eq:moving-parity}
\end{equation}
Its weighted size and the weighted first-alias defect share the growth
exponent
\begin{equation}
 \chi=\frac{\log(\beta R)}{\log\lambda}
 =0.463406944517002\ldots .
 \label{eq:common-exponent}
\end{equation}
\end{theorem}

\begin{proof}
Taylor's theorem for $f(\delta)=1-(1-\delta)^n$ gives
\[
 f(\delta)=n\delta-\frac{n(n-1)}2(1-\xi)^{n-2}\delta^2
\]
for some $0\le\xi\le\delta$, proving \eqref{eq:uniform-remainder}.
On \eqref{eq:natural-clock}, $k_\sigma\delta_\sigma\to0$, so the remainder
is $o(k_\sigma\delta_\sigma r_H^{-2k_\sigma})$ and
\eqref{eq:moving-parity} follows.

The archived multiplier law is
\begin{equation}
 \beta_k=\beta\exp\left[-\frac{\log C_M}{2k}+o(k^{-1})\right],
 \qquad C_M=1.946342905200968\ldots .
 \label{eq:beta-k}
\end{equation}
Hence
\begin{equation}
 \frac{\cA_{k,2k}}{2k}
 =\frac{\beta^{2k}}{C_M}\{1+o(1)\}.
 \label{eq:alias-asymptotic}
\end{equation}
After multiplication by $R^{2k}$, this is a constant multiple of
$(\beta R)^{2k}$.  The weighted parity packet divided by $2k$ is
$C_*\sqrt\sigma(R/r_H)^{2k}\{1+o(1)\}$.  Using
$2k=\log(1/\sigma)/\log\lambda+O(1)$ gives the same exponent
\eqref{eq:common-exponent}.
\end{proof}

\section{Clock phase and a scalar-only obstruction}

The bounded integer phase must be retained:
\begin{equation}
 \eta_\sigma
 =k_\sigma-\frac{\log(1/\sigma)}{2\log\lambda}.
 \label{eq:phase}
\end{equation}

\begin{theorem}[Phase law and unique symbolic scalar balance]
\label{thm:phase-law}
Along every natural-clock sequence with bounded $\eta_\sigma$,
\begin{equation}
 \boxed{
 \frac{\cP_{\sigma,2k_\sigma}}{\cA_{k_\sigma,2k_\sigma}}
 =C_*C_M\lambda^{\eta_\sigma}\{1+o(1)\}.}
 \label{eq:ratio-law}
\end{equation}
Consequently, along a subsequence with $\eta_\sigma\to\eta$, the ratio tends
to $C_*C_M\lambda^\eta$.  The leading scalar packets balance at the unique
symbolic phase
\begin{equation}
 \eta_*=-\frac{\log(C_*C_M)}{\log\lambda}.
 \label{eq:balance-phase}
\end{equation}
\end{theorem}

\begin{proof}
Divide \eqref{eq:moving-parity} by \eqref{eq:alias-asymptotic} and use
$\beta^{2k}=r_H^{-2k}\lambda^{-k}$:
\[
 \frac{\cP_{\sigma,2k}}{\cA_{k,2k}}
 =C_*C_M\sqrt\sigma\,\lambda^k\{1+o(1)\}.
\]
Equation \eqref{eq:phase} makes
$\sqrt\sigma\lambda^k=\lambda^{\eta_\sigma}$, proving
\eqref{eq:ratio-law}.  Positivity of $C_*,C_M$, and $\lambda>1$ makes the
right side strictly increasing in $\eta$, so its value one has the unique
solution \eqref{eq:balance-phase}.
\end{proof}

\begin{corollary}[Conditional canonical-phase scalar obstruction]
\label{cor:conditional-obstruction}
Put
\begin{equation}
 q_+=C_*C_M\lambda.
 \label{eq:q-plus}
\end{equation}
Assume $q_+<1$, let $|\eta_\sigma|\le1$, and suppose
$\cT_{\sigma,2k}=o(\cA_{k,2k})$.  Then
\begin{equation}
 |e_{\sigma,k,2k}|
 \ge\{1-q_++o(1)\}\cA_{k,2k},
 \label{eq:scalar-obstruction}
\end{equation}
and this is not $o(kR^{-2k})$.
\end{corollary}

\begin{proof}
Equation \eqref{eq:ratio-law} and $|\eta_\sigma|\le1$ give
$\cP_{\sigma,2k}/\cA_{k,2k}\le q_++o(1)$.  Substitute into
\eqref{eq:packet-form}.  Finally,
\[
 \frac{\cA_{k,2k}}{kR^{-2k}}
 =\frac{2}{C_M}(\beta R)^{2k}\{1+o(1)\}\longrightarrow\infty
\]
because $\beta R>1$.
\end{proof}

This is deliberately a conditional scalar-only negative result.  The
repository proves the symbolic corollary but does not contain a directed
interval enclosure certifying $q_+<1$.  The archived decimal values make the
hypothesis numerically plausible; they are not used as a rigorous inequality.
Even under that hypothesis, the result does not bound the actual raw trace
packet and does not prove divergence of the combined physical residual.

The same phase controls the local clearance.  If the archived boundary-cycle
clearance is
\begin{equation}
 \Delta_k^{\rm clr}=C_{\rm b}\lambda^{-2k}\{1+o(1)\},
 \qquad C_{\rm b}>0,
 \label{eq:clearance}
\end{equation}
then
\begin{equation}
 d_{\sigma,k}:=\frac{\Delta_k^{\rm clr}}\sigma
 =C_{\rm b}\lambda^{-2\eta_\sigma}\{1+o(1)\}.
 \label{eq:clearance-phase}
\end{equation}
At scalar balance,
\begin{equation}
 d_*=C_{\rm b}(C_*C_M)^2.
 \label{eq:balance-clearance}
\end{equation}
Thus suppressing the clock phase would discard both the symbolic scalar
balance ratio and the RH-322 entrance profile.  The decimal evaluations of
$C_{\rm b}$ and $d_*$ are numerical diagnostics only.

\section{Typed interface for the neighboring shell}

RH-322--RH-324 retain the entrance clearance and the oriented affine frame
\begin{equation}
 V\sim g_d,
 \qquad U=-\alpha V-Z_1,
 \qquad W=-\lambda U+Z_2,
 \qquad (V,U,W)\text{ has orientation }(+,-,+),
 \label{eq:local-frame}
\end{equation}
where $\alpha=2u_c$ and the output is compared after the center shift
$\kappa_{\rm aff}d=2u_c\lambda d$.  Their objects are forward probability
laws of mass one.  They do not yet define a signed trace observation, and
RH-19 shows that the neighboring critical sibling cannot be hidden in a far
Gaussian tail.

The exact handoff is therefore conditional but typed.  A later theorem must
construct, on the same clock, phase, coordinate frame, and Hardy
normalization, a decomposition
\begin{equation}
 \cT_{\sigma,2k}
 =\cB_{\sigma,k}(d,\eta;V,U,W,\mathfrak t)
 +\cS_{\sigma,k}(d,\eta;\mathfrak t)
 +\cR_{\sigma,k}.
 \label{eq:future-decomposition}
\end{equation}
Here $\mathfrak t$ is the still-missing trace observation, $\cB$ is the
physical boundary packet, $\cS$ is the neighboring-shell packet, and $\cR$
contains the second-leg, all-leg transport, and observation remainders.
Combining \eqref{eq:future-decomposition} with
\cref{thm:main-identity} gives the exact join
\begin{equation}
 e_{\sigma,k,2k}
 =\cB_{\sigma,k}+\cS_{\sigma,k}+\cR_{\sigma,k}
 +\cP_{\sigma,2k}-\cA_{k,2k}.
 \label{eq:typed-join}
\end{equation}
Consequently, if $\cR_{\sigma,k}=o(kR^{-2k})$, the desired matching is
equivalent to the signed condition
\begin{equation}
 \boxed{
 \cB_{\sigma,k}+\cS_{\sigma,k}
 =\cA_{k,2k}-\cP_{\sigma,2k}+o(kR^{-2k}).}
 \label{eq:shell-target}
\end{equation}
Equation \eqref{eq:shell-target} is an interface, not a theorem that the
packets on its left have been constructed.

\section{Reproduction and claim boundary}

The reproduction code checks the exact moment split, the parity signs, the
uniform Taylor bound, the clock/clearance dictionary, and the scalar balance
law.  Ordinary floating-point substitution of the archived decimal constants
gives the following diagnostic ledger:
\begin{center}\small
\begin{tabular}{rrr}
\toprule
$\eta$ & $C_*C_M\lambda^\eta$ & $C_{\rm b}\lambda^{-2\eta}$\\
\midrule
$-1$ & $0.122049588$ & $1.298368749$\\
$0$ & $0.204869205$ & $0.460805149$\\
$1$ & $0.343888020$ & $0.163544745$\\
$3.060914914$ & $1.000000000$ & $0.019340633$\\
\bottomrule
\end{tabular}
\end{center}
These are evaluations of proved symbolic formulas, not fitted noisy traces or
directed-rounding enclosures.  In particular, the table does not certify
$q_+<1$, $\eta_*>1$, or the displayed clearance interval.

\begin{remark}[Claim boundary]
The exact identity combines only known algebraic packets.  The phase law is
symbolic, while the canonical-window obstruction is conditional on
$q_+<1$ because the archived decimals are not interval certified.  The paper
does not
identify the RH-322--RH-324 probability law with $\cT_{\sigma,2k}$, control
the second physical critical leg, prove all-leg phase transport, bound the
weighted trace observation, include the neighboring shell, or establish
\eqref{eq:future-decomposition}.  Hence no joint first-alias trace law,
full-trace replacement, or actual divergence theorem is obtained.  Gates
A--E remain false/open.  This paper constructs no Hilbert--Polya operator,
identifies no Riemann zero, proves no von Mangoldt trace formula or
zeta-divisor equality, and does not imply RH.
\end{remark}

RH-327 should now compute the neighboring-shell coupling in the typed frame
of \eqref{eq:future-decomposition}; RH-328 may formulate the joint matching
equation only after that shell and the trace observation are defined.

\end{document}
